% SOURCE_STATEMENT: "Determine the relation between the special boundary quantum Knizhnik–Zamolodchikov solutions obtained as quantum correlation functions or $q$-integrals and the affine-Hecke-algebra basis characterized by its asymptotics in a Weyl chamber."
% Preparation and final self-audit: CONTRIBUTING.md prompts C and D; source/audit artifacts retained in the local physics_statements_2026_10_11 batch.
% !TeX program = lualatex
\documentclass[11pt]{article}
\usepackage[margin=25mm]{geometry}
\usepackage{fontspec}
\setmainfont{Latin Modern Roman}
\usepackage{amsmath,amssymb,mathtools}
\usepackage{unicode-math}
\setmathfont{Latin Modern Math}
\usepackage{xurl,hyperref}
\hypersetup{colorlinks=true,urlcolor=blue}
\setcounter{secnumdepth}{0}
\setlength{\parindent}{0pt}
\setlength{\parskip}{5pt}
\setlength{\emergencystretch}{3em}
\newcommand{\sref}[2]{\hyperlink{source:#1:#2}{[S#2]}}
\newcommand{\cataloguescope}[1]{\paragraph{Recorded target.}#1}
\newcommand{\R}{\mathbb R}
\newcommand{\C}{\mathbb C}
\newcommand{\N}{\mathbb N}
\newcommand{\Z}{\mathbb Z}
\newcommand{\Q}{\mathbb Q}
\newcommand{\F}{\mathbb F}
\newcommand{\eps}{\varepsilon}
\begin{document}
\section{Quantum-Group and Affine-Hecke Solutions of Boundary qKZ Equations}\label{statement}

\subsection{Definitions and mathematical statement}

\paragraph{Shared notation.}$\mathbb N=\{1,2,\ldots\}$; the standard number fields have their usual meanings. All additional conventions are stated locally.


\paragraph{Definitions and conventions.}
The boundary quantum Knizhnik--Zamolodchikov (qKZ) equations are compatible difference equations for vector-valued functions of spectral parameters, with exchange operators built from $R$- and boundary $K$-matrices. An affine Hecke algebra is the deformation of an affine Weyl-group algebra subject to braid and quadratic relations. A Weyl chamber is a region determined by inequalities between root coordinates. A $q$-integral is the discrete analogue of integration adapted to multiplicative $q$-shifts.

\paragraph{Target retained from the canonical source catalogue.}
\paragraph{Catalogue wording retained for traceability.}
Determine the relation between the special boundary quantum Knizhnik–Zamolodchikov solutions obtained as quantum correlation functions or $q$-integrals and the affine-Hecke-algebra basis characterized by its asymptotics in a Weyl chamber.

\paragraph{Source scope and qualification.}
This comparison depends on matching parameters, normalizations and analytic domains of the two constructions; an arbitrary basis identification is not the requested relation.
\subsection{Short English statement}
Relate boundary qKZ solutions built by quantum groups or q-integrals to the affine-Hecke basis selected by chamber asymptotics.
\subsection{Sources}
\begin{itemize}
\item[\hypertarget{source:statement:1}{S1}] ULAM.AI and the UnsolvedMath Contributors, UnsolvedMath, catalogue entry. Catalogue curation under CC BY 4.0; definitions and English exposition added locally.
\url{https://huggingface.co/datasets/ulamai/UnsolvedMath}.
\item[\hypertarget{source:statement:2}{S2}] Mini-Workshop: Coideal Subalgebras of Quantum Groups (2015), Oberwolfach Reports, DOI 10.4171/owr/2015/10, pp. 533--569. Primary report accessed 2026-10-11; relevant PDF pages 31.
\url{https://doi.org/10.4171/owr/2015/10}.
\end{itemize}
\label{end:statement}
\end{document}
